\chapter{Dati e configurazione}
\label{chap:dati}

\section{Organizzazione di codice e dati}

Le immagini non devono essere copiate nel repository. Per una lavorazione
industriale è più ordinato mantenere una directory esterna per dataset e
risultati, usando \filepath{.env} per collegarla al codice. Il checkpoint YOLO
è invece conservato nella directory \filepath{models/} del progetto. Una
possibile organizzazione dei dati esterni è:

\begin{lstlisting}[language={}]
D:/GraphReconstructionData/
  datasets/
    impianto_A/
      images/
      labels_extended/
  work/
    impianto_A/
\end{lstlisting}

La separazione limita la commistione tra sorgenti, dati sensibili e cache e
consente di conservare più lavorazioni indipendenti.

\subsection{Convenzione di denominazione delle immagini}

La discovery non visita sottocartelle e legge soltanto file con estensione
minuscola \code{.jpg} presenti direttamente in \envvar{IMAGES_DIR}. Il nome
deve rispettare la forma:

\begin{center}
  \code{node<id>_<frame_numerico>.jpg}
\end{center}

L'ultima parte prima dell'estensione è sempre il numero del frame. Tutto ciò che
si trova tra \code{node} e l'ultimo underscore è l'ID della lavorazione. Sono
quindi validi entrambi questi gruppi:

\begin{lstlisting}[language={}]
node7_0001.jpg
node7_0002.jpg

node15_1_0001.jpg
node15_1_0002.jpg
\end{lstlisting}

File come \code{node7_frontale.jpg}, immagini PNG o estensioni \code{.JPG} non
vengono rilevati. Se un nodo non viene trovato, verificare prima nomi ed
estensioni.

\section{Configurazione locale}

Nella root del progetto è presente \filepath{.env.template}. Copiarlo in
\filepath{.env}; il nuovo file resta locale e non deve contenere dati da
versionare:

\begin{lstlisting}[language={}]
Copy-Item .env.template .env
\end{lstlisting}

Per l'esempio precedente, la parte principale del file diventa:

\begin{lstlisting}[language={}]
IMAGES_DIR=D:/GraphReconstructionData/datasets/impianto_A/images
LABELS_EXT_DIR=D:/GraphReconstructionData/datasets/impianto_A/labels_extended
YOLO_MODEL_PATH=./models/yolo_detector.pt

ASSETS_DIR=./assets
ANNOTATIONS_DIR=./annotations
WORK_DIR=D:/GraphReconstructionData/work/impianto_A
\end{lstlisting}

Le prime tre variabili sono obbligatorie. \envvar{IMAGES_DIR} e
\envvar{YOLO_MODEL_PATH} servono alla ricostruzione;
\envvar{LABELS_EXT_DIR} è usata soprattutto da \filepath{evaluate.py}. In
assenza di label può indicare una directory vuota; la valutazione supervisionata
verrà omessa.

\envvar{WORK_DIR} è la directory radice degli output: la pipeline vi crea le
cache delle detection, i risultati multivista, i confronti e gli eventuali file
diagnostici. \envvar{ASSETS_DIR} e \envvar{ANNOTATIONS_DIR} indicano invece le
risorse di supporto descritte di seguito.

I percorsi relativi, come \filepath{./assets}, sono risolti rispetto alla root
del repository, non rispetto alla directory dalla quale è stato aperto il
terminale. Per dataset e modelli esterni sono preferibili percorsi assoluti.

Se una variabile omonima esiste nell'ambiente di Windows o nella sessione
Conda, ha precedenza sul valore scritto in \filepath{.env}. In caso di percorsi
inattesi, verificare sia i valori caricati sia le variabili della sessione.

\subsection{Dati di riferimento}

La directory indicata da \envvar{ASSETS_DIR} contiene mapping, porte e primitive
SVG dei componenti Shurikode. Queste risorse servono a rappresentare
correttamente i componenti nelle visualizzazioni CAD e QGIS, ma non modificano
la topologia del grafo ricostruito. Le primitive attuali sono risorse di
progetto e potranno essere aggiornate o adattate agli standard aziendali.

La directory indicata da \envvar{ANNOTATIONS_DIR} contiene i grafi annotati con
nome \code{node<id>_gt.json}. Questi file non servono a produrre
\filepath{fused_graph.json}; servono a confrontare il risultato con un grafo
atteso. Anche i grafi campione sotto
\filepath{evaluation/standard_comparison/} appartengono alla fase di
valutazione, non alla ricostruzione.

\Needspace{17\baselineskip}
\subsection{Formato delle annotazioni di grafo}

Ogni file di annotazione descrive i componenti di un nodo idraulico e le
connessioni attese tra essi. Il formato è il seguente:

\begin{lstlisting}[language={}]
{
  "node_id":"11", "best_view":"node11_0003.jpg",
  "nodes": [
    {"nid":"95#1","code":95,"name":"Bout bic.","valence":2,"cx":0.58,"cy":0.31},
    {"nid":"15#2","code":15,"name":"Saracinesca","valence":2,"cx":0.57,"cy":0.34}
  ],
  "edges": [["95#1", "15#2"]]
}
\end{lstlisting}

\code{node_id} identifica il gruppo di immagini e \code{best_view} la vista
usata come riferimento per l'annotazione. In \code{nodes}, \code{nid} è
l'identificatore univoco nel formato \code{<code>\#<instance_id>}; \code{code}
è il codice Shurikode, \code{valence} il numero atteso di connessioni e
\code{cx}, \code{cy} le coordinate normalizzate del centro. \code{edges}
contiene infine le coppie di \code{nid} collegate nel grafo atteso.

\subsection{Opzioni VGGT e debug}

Per il primo run si raccomanda di mantenere i valori VGGT del template. In particolare,
\envvar{VGGT_EDGE_ENABLED=true} abilita la geometria 3D e
\envvar{VGGT_GENERATE_IF_MISSING=true} consente di generare gli artefatti
mancanti. La prima elaborazione può quindi essere sensibilmente più lenta delle
successive.

Il debug dettagliato va inizialmente lasciato spento:

\begin{lstlisting}[language={}]
DEBUG_OUTPUT_ENABLED=false
DEBUG_OUTPUT_DIR=./output/debug
\end{lstlisting}

Il debug va abilitato soltanto per l'analisi di nodi problematici.

\section{Verifica della configurazione}

Con l'ambiente Conda attivo, si può verificare che la configurazione sia
caricabile senza avviare YOLO:

\begin{lstlisting}[language={}]
python -c "from graph_reconstruction.config import IMAGES_DIR, WORK_DIR; print(IMAGES_DIR); print(WORK_DIR)"
\end{lstlisting}

I due percorsi stampati devono corrispondere alla lavorazione corrente. Se il
comando segnala \code{Missing required path setting}, il file \filepath{.env}
non esiste, non è nella root corretta oppure manca una delle variabili
obbligatorie.
